\capitulocapa{images/cap-5.png}{Conhecendo o Zustand}
\label{chap:conhecendo-zustand}

\section{A terceira alteração no mesmo Provider}

\begin{historia}
Depois de separar o Context no capítulo anterior, o tema estava organizado. O
Provider era pequeno, os consumidores eram conhecidos e a aplicação funcionava.

Então a equipe recebeu outra necessidade transversal: o Drawer de notificações
deveria abrir pelo Header, pelo Dashboard e por atalhos de teclado.

Alguém sugeriu criar outro Context. Outra pessoa perguntou quantos Providers
teríamos daqui a seis meses. Foi quando apareceu uma alternativa:

``E se esse estado não precisasse morar dentro da árvore de componentes?''

A pergunta levou o time ao Zustand.
\end{historia}

\ilustracaopendente{5-1-terceira-alteracao}{Cena em um ambiente de trabalho:
três pessoas observam o DevPanel em um monitor. Na tela, o Header, o Dashboard e
um atalho de teclado apontam para o mesmo Drawer de notificações. Ao lado do
monitor, vários blocos identificados como ``Provider'' começam a se empilhar,
enquanto uma das pessoas sugere retirar esse estado da árvore de componentes.}

Zustand não entra no DevPanel porque Context falhou. Context continua apropriada
para vários cenários. Zustand entra porque começamos a precisar de stores com
acesso direto, assinaturas mais específicas e uma organização que não dependa da
posição dos componentes.

\section{O problema: distribuição e gerenciamento se misturaram}

Context distribui um valor produzido por algum componente. Por isso, normalmente
criamos um Provider, colocamos estado dentro dele, definimos ações e envolvemos
uma parte da árvore.

Essa solução une duas responsabilidades:

\begin{itemize}
  \item o componente Provider participa do ciclo de vida do React;
  \item o contexto transporta o valor até os consumidores.
\end{itemize}

Para o tema, o custo era aceitável. Conforme novos estados transversais aparecem,
começamos a repetir Providers, hooks de validação, memoização e decisões sobre
onde envolver a árvore.

O que procuramos agora é uma fonte de verdade que exista fora dos componentes e
permita que cada consumidor assine apenas aquilo que utiliza.

\section{Como normalmente resolvemos}

Antes de adicionar uma dependência, vale lembrar as opções que já conhecemos.

\subsection{Elevar o estado e passar props}

É a solução mais explícita. Funciona muito bem quando os consumidores estão
próximos e a árvore é estável. O custo surge com intermediários e rotas distantes.

\subsection{Criar um Context pequeno}

É adequado para valores transversais e pouco mutáveis. O custo aparece quando
precisamos de várias atualizações independentes, seleção granular ou acesso fora
de componentes.

\subsection{Usar uma variável de módulo}

Uma variável exportada existe fora da árvore, mas não é uma store React. Alterar
o valor não cria por si só uma assinatura, não solicita renderização e não
estabelece um contrato seguro de atualização.

Zustand ocupa justamente esse espaço: um store externo com mecanismo de
assinatura integrado ao React.

\section{O que é uma store do Zustand?}

Uma store reúne estado e ações. Ela existe independentemente de qualquer
componente. O React se conecta a ela por um hook, e cada componente informa qual
parcela deseja observar.

Uma primeira store de tema pode ser escrita assim:

\begin{lstlisting}[caption={Uma store pequena de tema.},label={lst:theme-store}]
import { create } from 'zustand'

type Theme = 'light' | 'dark'

type ThemeStore = {
  theme: Theme
  toggleTheme: () => void
}

export const useThemeStore = create<ThemeStore>((set) => ({
  theme: 'light',
  toggleTheme: () =>
    set((state) => ({
      theme: state.theme === 'light' ? 'dark' : 'light',
    })),
}))
\end{lstlisting}

Mesmo antes de estudar cada detalhe, conseguimos enxergar o contrato:
\codigo{theme} é o dado e \codigo{toggleTheme} é a operação permitida. Não há
Provider nem prop intermediária.

\begin{decisao}
Uma store não deve ser chamada de global apenas porque pode ser importada em
qualquer lugar. Seu alcance técnico é amplo; sua responsabilidade deve continuar
pequena e explícita.
\end{decisao}

\section{Consumindo apenas o necessário}

O Header precisa do valor e da ação:

\begin{lstlisting}
export function Header() {
  const theme = useThemeStore((store) => store.theme)
  const toggleTheme = useThemeStore((store) => store.toggleTheme)

  return (
    <header>
      <span>Tema: {theme}</span>
      <button type="button" onClick={toggleTheme}>
        Alternar tema
      </button>
    </header>
  )
}
\end{lstlisting}

A função entregue ao hook é um \emph{selector}. Ela declara qual parte interessa
ao componente. Quando \codigo{theme} mudar, o Header recebe o novo valor. Quando
outra informação da store mudar no futuro, esse selector ajuda a evitar uma
atualização sem relação com o Header.

A Sidebar precisa apenas ler:

\begin{lstlisting}
export function Sidebar() {
  const theme = useThemeStore((store) => store.theme)

  return <aside>Tema atual: {theme}</aside>
}
\end{lstlisting}

Não passamos \codigo{theme} pelo Layout. Não envolvemos a aplicação num Provider.
Header e Sidebar importam o mesmo hook, mas cada um possui sua própria assinatura.

\section{O que acontece quando a ação é chamada?}

A função \codigo{create} recebe uma criadora de estado. Zustand fornece
\codigo{set}, usado para produzir a próxima versão da store.

No nosso exemplo, o próximo tema depende do atual. Por isso usamos a forma
funcional:

\begin{lstlisting}
set((state) => ({
  theme: state.theme === 'light' ? 'dark' : 'light',
}))
\end{lstlisting}

Quando a ação executa:

\begin{enumerate}
  \item Zustand lê o estado atual;
  \item a função devolve a parcela que mudou;
  \item a store combina a atualização no nível superior;
  \item as assinaturas são notificadas;
  \item consumidores cujo resultado selecionado mudou renderizam novamente.
\end{enumerate}

Por baixo da integração com React existe o protocolo apropriado para stores
externos. Não precisamos manipular o DOM, emitir eventos próprios ou implementar
uma inscrição manual.

\section{Antes e depois}

Na versão com props, o Layout possuía o tema e o entregava aos filhos. Na versão
com Context, um Provider mantinha o valor e o distribuía. Com Zustand, a store é
a fonte de verdade, e os componentes assinam diretamente suas parcelas.

Isso reduz código de transporte. Não elimina a necessidade de arquitetura.
Ainda precisamos decidir:

\begin{itemize}
  \item o que merece entrar na store;
  \item quais ações fazem parte de seu contrato;
  \item quanto tempo o estado deve existir;
  \item se o dado pertence ao cliente ou ao servidor;
  \item quais componentes realmente precisam assiná-lo.
\end{itemize}

Zustand facilita a mecânica. As decisões continuam sendo nossas.

\ilustracaopendente{5-2-props-context-zustand}{Diagrama em três quadros. No
primeiro, o estado desce pela árvore por várias props intermediárias. No segundo,
um Provider envolve a árvore e distribui o valor por Context. No terceiro, uma
store do Zustand fica ao lado da árvore, e Header e Sidebar se conectam
diretamente a parcelas específicas por meio de selectors. Destacar visualmente
que somente os consumidores interessados recebem cada atualização.}

\section{Limitações e cuidados}

\subsection{Importar de qualquer lugar é tentador}

Como a store não depende da árvore, qualquer arquivo pode importá-la. Se o time
usar essa liberdade sem fronteiras, dependências invisíveis se espalham. Uma
store pequena deve possuir nome, domínio e contrato reconhecíveis.

\subsection{Nem todo useState deve migrar}

O campo de busca usado por uma única página continua local. A abertura de um
menu pertencente a um componente continua local. Mover tudo para Zustand aumenta
o alcance e prolonga o ciclo de vida sem necessidade.

\subsection{Dados remotos possuem outro problema}

A lista de usuários vem de uma API. Ela envolve cache, atualização, erro e
sincronização com o servidor. Conseguimos colocá-la numa store, mas isso não
significa que deveríamos reconstruir manualmente uma ferramenta de consultas.
Esse será o assunto do próximo capítulo.

\section{Quando considerar Zustand}

Zustand começa a fazer sentido quando temos estado do cliente compartilhado por
regiões distantes, quando selectors reduzem o acoplamento das atualizações ou
quando a lógica precisa ser acessada fora dos componentes.

Também é útil quando a equipe quer ações próximas do estado sem adotar uma
estrutura muito extensa. A API pequena é uma vantagem desde que não seja usada
como desculpa para uma store sem organização.

Não existe uma quantidade mágica de componentes. O critério é o custo que o
compartilhamento atual impõe e a natureza do dado.

\section{Exercício: migre o tema}

\begin{tarefa}
Substitua \codigo{ThemeProvider} por uma store de tema. Conecte Header, Sidebar e
Perfil diretamente ao hook da store. Remova o Provider e todas as props de tema
que restarem na árvore.
\end{tarefa}

Depois da migração, registre:

\begin{enumerate}
  \item onde ficou a fonte de verdade;
  \item quais componentes assinam o valor;
  \item quais componentes assinam a ação;
  \item quais arquivos deixaram de transportar propriedades;
  \item quais estados do DevPanel devem continuar locais.
\end{enumerate}

\espacoresposta{9}

Evite adicionar filtros, sessão ou dados da API à store neste exercício. Uma
primeira store pequena ensina mais que uma store ``completa'' sem fronteiras.

\begin{resumo}
\begin{itemize}
  \item Zustand cria stores externos à árvore de componentes.
  \item Estado e ações formam o contrato público da store.
  \item Componentes usam selectors para assinar parcelas específicas.
  \item O acesso direto elimina props de transporte e a exigência de Provider
    no uso comum.
  \item A API pequena não elimina decisões sobre responsabilidade e ciclo de
    vida.
  \item Estado local e Context continuam úteis quando combinam melhor com o
    problema.
\end{itemize}
\end{resumo}

Conhecer Zustand não significa colocar tudo nele. Significa ter outra ferramenta
e saber reconhecer o problema para o qual ela foi desenhada.
